\documentclass{article}

\usepackage{tabularx}
\usepackage{booktabs}
\usepackage{longtable}
\usepackage{float}
% Extra vertical space between paragraphs (no first-line indent) for readability.
\usepackage[skip=0.55\baselineskip plus 2pt minus 1pt]{parskip}

\title{Reflection and Traceability Report on \progname}

\author{\authname}

\date{}

\input{../Comments}
\input{../Common}

\begin{document}

\maketitle

%\plt{Reflection is an important component of getting the full benefits from a
%learning experience.  Besides the intrinsic benefits of reflection, this
%document will be used to help the TAs grade how well your team responded to
%feedback.  Therefore, traceability between Revision 0 and Revision 1 is and
%important part of the reflection exercise.  In addition, several CEAB (Canadian
%Engineering Accreditation Board) Learning Outcomes (LOs) will be assessed based
%on your reflections.}

\section{Changes in Response to Feedback}

\noindent
This section summarizes feedback from TAs and course rubrics, instructor expectations,
peer and design reviews, teammates, cross-team coordination (where applicable), the
project supervisor (including post--Rev~0 alignment after the standalone MacSync pivot---see
the Reflection on Project Management section for meeting frequency), and structured
usability-test participants.
Each item below states the \emph{source}, the \emph{feedback}, and the \emph{change}
(documentation, process, or code) so graders can trace rationale without opening every link.\\

% TA/Rubric Feedback Traceability: Added explicit subsections for Problem Statement/Goals and Dev Plan so feedback-to-change mapping is easy to grade and does not get lost under unrelated deliverables.
\subsection{Problem Statement and Goals}
The table below summarizes the feedback and the changes made in response to it.
% TA/Rubric Feedback Traceability: Summarized each rubric comment and the corresponding document change in a compact table, per the template's guidance to address every feedback item.
\begin{table}[H]
\caption{Problem Statement and Goals --- Feedback Response Summary} \label{TblPSGFeedbackResponse}
\begin{tabularx}{\textwidth}{lXX}
\toprule
\textbf{Source} & \textbf{Feedback Item (Brief)} & \textbf{Change Made (Brief)} \\
\midrule
TA Rubric & Main objective/focus unclear; integrated project context missing & Added explicit project objective, scope focus, and integration context in problem statement. \\
TA Rubric & Missing discussion: security/privacy, platform scope, integrations, performance & Added security/privacy concerns, mobile+web scope, third-party integration notes, and high-traffic expectations. \\
TA Rubric & Stakeholder hierarchy too rigid/over-structured & Replaced primary/secondary/tertiary hierarchy with a single stakeholder list reflecting overlaps. \\
TA Rubric & Goals not clearly measurable; missing security/privacy goal & Rewrote goals with intent + measurable criteria + rationale; added Security \& Privacy goal. \\
TA Rubric & Reflection should be individual (Q1/Q2), team-wide (Q3) & Restructured reflection appendix: Q1/Q2 per member; Q3 team response. \\
\bottomrule
\end{tabularx}
\end{table}

% TA/Rubric Feedback Traceability: Added a concise summary subsection for Development Plan feedback so Rev 1 changes can be traced without requiring the reader to infer them from the full document.
\subsection{Development Plan and Proof of Concept Demonstration Plan}
The table below summarizes the feedback and the changes made in response to it.
% TA/Rubric Feedback Traceability: Captured both the rubric's content-depth comments and the LO13 tool-selection justification feedback in one place.
\begin{table}[H]
\caption{Development Plan --- Feedback Response Summary} \label{TblDevPlanFeedbackResponse}
\begin{tabularx}{\textwidth}{lXX}
\toprule
\textbf{Source} & \textbf{Feedback Item (Brief)} & \textbf{Change Made (Brief)} \\
\midrule
TA Rubric (LO13) & Section 10 tech choices lacked rationale/pros/cons & Expanded ``Expected Technology'' with project-specific rationale, trade-offs, and limitations per tool/tech. \\
TA Rubric & Roles missing technical responsibilities and overlap plan & Added technical ownership (primary/secondary) across frontend, mobile, backend, DB, CI/CD. \\
TA Rubric & Workflow not at Contributor's-Guide level & Added PR/review expectations, labels, quality gates, and Definition of Done. \\
TA Rubric & Scheduling section shallow; requested timetable-style detail & Added a high-level timetable table with phases and planned outcomes. \\
TA Rubric & POC plan needed more context + failed-POC consequences & Added project context, expanded risks, added scenario-based demo + failure-mode check, defined failed POC. \\
TA Rubric & Formatting: justified alignment not applied in places & Removed \texttt{flushleft} usage in communication plan section to avoid forced left alignment. \\
TA Rubric & GitHub link should point to main repo page & Updated link from board to repository URL. \\
TA Rubric & Reflection should be individual (Q1/Q2), team-wide (Q3) & Restructured reflection appendix: Q1/Q2 per member; Q3 team response. \\
\bottomrule
\end{tabularx}
\end{table}

\subsection{SRS and Hazard Analysis}

\noindent\textbf{Software Requirements Specification (SRS).}\\
Rev~1 SRS content lives in the AsciiDoc sources under \path{docs/SRS/} (built to \path{docs/SRS/index.pdf}).
Table~\ref{TblSRSFeedbackResponse} summarizes feedback and edits; deliverable-level Rev~0$\rightarrow$~Rev~1
summary also appears in Table~\ref{TblRev0Rev1Trace}.

\begin{table}[H]
\caption{SRS --- Feedback Response Summary} \label{TblSRSFeedbackResponse}
\begin{tabularx}{\textwidth}{lXX}
\toprule
\textbf{Source} & \textbf{Feedback Item (Brief)} & \textbf{Change Made (Brief)} \\
\midrule
TA / peer review & Some functional and NFR statements were hard to verify directly from acceptance-style tests & Clarified expected behaviour, error paths, and measurable NFR targets; aligned requirement wording with derivations in \path{docs/VnVPlan/VnVPlan.tex} and results in \path{docs/VnVReport/VnVReport.tex}. \\
TA Rubric & Traceability between requirements and verification artifacts needed to be obvious for grading & Kept stable FR/NFR identifiers across SRS parts, MIS/MG, VnV Plan, and VnV Report; cross-checked module and scenario coverage against SRS sections (system goals and functional areas). \\
TA / readability & System context, scope, and naming should match the Hazard Analysis boundary and design docs & Harmonized LEMS scope, integration boundaries, and component terminology with \path{docs/HazardAnalysis/HazardAnalysis.tex} and the Module Guide; reduced cross-document naming drift. \\
Usability \& product feedback & Refund request and organizer event-edit flows were missing or place-held in the product; requirements needed to catch up & Updated scenario and requirement coverage for end-to-end refunds and working admin event edit, consistent with Rev~1 implementation and VnV evidence. \\
\bottomrule
\end{tabularx}
\end{table}

\noindent\textbf{Traceability between SRS and Hazard Analysis (bidirectional).}\\
The SRS specifies behaviours and qualities (FR/NFR) for LEMS/MacSync; the Hazard Analysis
(\path{docs/HazardAnalysis/HazardAnalysis.tex}) records FMEA hazards on the same subsystems (payments,
access control, sign-ups) and derives SSR-* safety and security requirements with concrete mitigations.
\textbf{SRS $\rightarrow$ HA:} SRS-scoped workflows and data (e.g.\ payments, RBAC/FBAC, capacity, notifications)
are the same capabilities the HA analyses for failure and abuse; HA assumptions and the system-boundary
figure align with the integrations and actors described in SRS environment and system parts.
\textbf{HA $\rightarrow$ SRS:} SSR requirements and roadmap entries in HA are reflected in SRS non-functional,
security, and operational expectations that VnV tests cite alongside FR/NFR IDs (system tests in the VnV Plan
often name both a functional requirement and a hazard or mitigation identifier, e.g.\ payment and breach scenarios).
Together, the SRS + HA pair explains \emph{what} must hold and \emph{what can go wrong}, with VnV linking both to tests.

\noindent\textbf{Hazard Analysis.}\\
Table~\ref{TblHazardAnalysisFeedbackResponse} summarizes TA feedback on the Hazard Analysis deliverable and the
corresponding edits in \path{docs/HazardAnalysis/HazardAnalysis.tex}.
% Traceability: maps TA HA feedback (boundary diagram + readability) to concrete document edits for Reflect \& Trace.
\begin{table}[H]
\caption{Hazard Analysis --- Feedback Response Summary} \label{TblHazardAnalysisFeedbackResponse}
\begin{tabularx}{\textwidth}{lXX}
\toprule
\textbf{Source} & \textbf{Feedback Item (Brief)} & \textbf{Change Made (Brief)} \\
\midrule
TA (overall) & A diagram in the section defining system boundaries would help show how components interact & Added a system-boundary diagram (boxed layout: LEMS web/mobile application areas, PostgreSQL, payment and notification integration interfaces to external providers) in \path{docs/HazardAnalysis/HazardAnalysis.tex} under \emph{System Boundaries and Components}, with caption and PDF-traceable figure label \texttt{fig:system-boundary}. \\
TA / readability & Boundary narrative should be easy to reconcile with the architecture and naming in the rest of the HA doc & Renamed boundary list wording to LEMS where inconsistent; rewrote the boundary paragraph to state in-scope components and that external payment/notification providers are treated as critical dependencies within the analysis boundary. \\
TA / rubric polish & Spelling, grammar, and professional presentation (supporting boundary and FMEA readability) & Removed a stray punctuation typo in the introduction workflow sentence; fixed ``dependecies,'' incorrect plural ``API's,'' and tightened out-of-scope API wording; corrected duplicated FMEA subfigure caption (Part~2). \\
\bottomrule
\end{tabularx}
\end{table}

\noindent\textit{Note:} Rubric scores for method depth, hazard completeness, and peer-review issue quality reflect process and prior drafts; the changes above address the documented TA request for a boundary diagram and alignment/readability of the boundaries section.

\subsection{Design and Design Documentation}

Design feedback is summarized in Table~\ref{TblDesignFeedbackResponse}; the narrative below adds file-level detail
and cross-team context.

\begin{table}[H]
\caption{Design documentation (MIS / MG) --- Feedback Response Summary} \label{TblDesignFeedbackResponse}
\begin{tabularx}{\textwidth}{lXX}
\toprule
\textbf{Source} & \textbf{Feedback Item (Brief)} & \textbf{Change Made (Brief)} \\
\midrule
Peer design review & MIS symbols and acronyms pointed readers to the SRS or MG instead of standing alone & Expanded in-document symbol/acronym table and added a per-module ``sufficient to build'' note in \path{docs/Design/SoftDetailedDes/MIS.tex}. \\
TA / team review & Module state for payments, sign-up, and access control was too implicit for an independent implementer & Rewrote M1/M3/M5 state with explicit structures (\texttt{SignUpRecords}, versioned policy storage, payment/refund records, \texttt{roleAssignments}); moved session-timeout expectations into the MIS. \\
Peer review & MG UI figures lacked orientation; abbreviations were inconsistent with MIS self-containment & Added explanatory text before figures, fixed a caption typo, and defined abbreviations inside \path{docs/Design/SoftArchitecture/MG.tex}. \\
Cross-team (History Manager) & External History Manager work needed a traceable, citable integration record & Documented Team~10 contribution with public PR \url{https://github.com/OmarHassanAdelhamid/Five-of-a-Kind-capstone-project-/pull/328}; completed MIS appendix placeholders tied to that collaboration. \\
\bottomrule
\end{tabularx}
\end{table}

\noindent
\textbf{Module Interface Specification (MIS).}\\
In response to design-review feedback, we strengthened the MIS so it better supports an independent implementer
 and stands alone as a document: the symbol and acronym table was expanded and explicitly scoped to this deliverable
  (no deferring readers to the SRS or MG for definitions).
We added a short subsection stating how each module section is intended to be sufficient to build from
 (interfaces, typed state, assumptions, behaviour).
Payment (M1), sign-up (M3), and access control (M5) state was rewritten using explicit data structures---for
 example consolidated payment and refund records, a named \texttt{SignUpRecords} map, and versioned policy
  storage (\texttt{policyByVersion}, \texttt{currentPolicyVersion}) with correct \texttt{roleAssignments} semantics.
Other edits included tightening the module hierarchy wording, stating session-timeout expectations inside
 the MIS instead of pointing to other docs, replacing appendix placeholder text, and documenting the History
  Manager work for Team~10 with a public pull request: \url{https://github.com/OmarHassanAdelhamid/Five-of-a-Kind-capstone-project-/pull/328}.
Unused SRS cross-document label imports were removed from the MIS source.\\

\noindent
\textbf{Module Guide (MG).}\\
We added brief narrative before each user-interface figure to explain what the screenshot shows and which modules it
 relates to, fixed a caption typo, and noted that abbreviations in the MG are defined in that document itself for
  consistency with the MIS approach.

\subsection{VnV Plan and Report}

\noindent\textbf{VnV Report:}
\begin{itemize}
\item \textbf{Traceability readability feedback:} FR and module traceability matrices were hard to read when scaled. \textbf{Change made:} moved matrices to become larger and wide matrices to landscape pages instead of shrinking width\texttt{\textbackslash resizebox}, improving readability without reducing font size.
\item \textbf{NFR consistency feedback:} report NFR labels did not consistently match the VnV Plan. \textbf{Change made:} aligned the report to plan IDs NFR1--NFR6 and added a dedicated NFR trace crosswalk table for one-to-one mapping.
\item \textbf{Scope/rigour feedback:} earlier wording implied full equivalence to planned load/usability activities. \textbf{Change made:} explicitly documented where lighter methods were used, what was partially verified, and what remains as scope gaps (for example webhook idempotency, reliability monitoring, and audit depth).
\item \textbf{Rev 0 product feedback:} ``request refund not yet working'' and ``admin cannot edit created event.'' \textbf{Change made:} implemented backend refund endpoint and frontend refund workflow, and replaced admin event-edit placeholder with a working edit form.
\end{itemize}

\noindent\textbf{VnV Plan:}\\
In response to TA feedback, we expanded the symbols and acronyms table
(beyond the placeholder row), removed instructor template prompts (\texttt{wss}) and replaced them
with final text, and tightened vague objectives and usability language---including concrete definitions
for elevated load (tied to existing NFR concurrency cases) and measurable usability criteria (task
success, timings, Likert item for NFR3). The appendix now documents symbolic parameters from those
tests and states how optional usability surveys relate to NFR3. After the Module Interface
Specification (MIS) was available, we filled the unit-test section: implementation verification
references modules M1--M8; for each module we added representative unit test cases (\texttt{UT-M*})
mapped to MIS access programs and exceptions, plus a traceability table to SRS areas.
Module-level cases live under the section titled \texttt{Unit Test Description} in \path{docs/VnVPlan/VnVPlan.tex}.
Revision history records version~1.1 (MIS-dependent content) and~1.2 (TA feedback pass).\\

\subsection{Implementation, coding standards, and repository completeness}

\noindent
\textbf{Code and standards.}\\
The team followed a consistent stack-idiomatic style (e.g.\ TypeScript/ESLint in frontend and backend packages)
with meaningful identifiers, module boundaries aligned to MIS modules M1--M8 in design reviews, and review gates for
payment and auth changes. TypeScript and shared enums reduce ambiguous numeric codes where practical;
reviewers are pointed to MIS access programs and exceptions as the contract for unit tests (\texttt{UT-M*}).\\

\noindent
\textbf{Repository completeness (capstone template expectations).}\\
Root-level onboarding files are present and maintained relative to the course template: \path{README.md},
\path{INSTALL.md}, \path{CONTRIBUTING.md}, and \path{LICENSE}. Per-component instructions exist under
\path{src/frontend/README.md}, \path{src/backend/README.md}, and selected \path{docs/*/README.md} folders.
Continuous integration runs automated checks (e.g.\ GitHub Actions) on pull requests to catch regressions before merge.
This reflection does not replace line-level code review; it documents where graders should look for style, CI, and
module-to-code mapping alongside the MIS and repo layout.

\section{Extras}

\noindent\textit{Grading note:} this section is split into \textbf{Extra~1} and \textbf{Extra~2} so rubric lines map cleanly to identifiable deliverables.

\subsubsection*{Extra 1 --- User documentation (Member and Staff/Admin guides)}

This project completed instructor-approved extra user documentation aligned with the capstone problem-statement
goals for an MES-facing event platform (member mobile experience and staff administration).

We produced two companion PDF manuals under
\url{docs/Extras/UserManual/},
each version~1.0.0 (April~2026):
\begin{itemize}
  \item \textbf{MacSync Member Guide} (\path{User_Manual.tex}) --- end-user
        documentation for McMaster students using the iOS/Android app: account
        creation and \texttt{@mcmaster.ca} verification, navigation (Events, My
        Tickets, Notifications, Profile), registering for free and paid events
        (including Stripe checkout and optional table/bus seat selection), viewing
        tickets and entry QR codes, requesting refunds, notification types, and
        profile settings with a short glossary.
  \item \textbf{MacSync Staff \& Admin Guide} (\path{Admin_Manual.tex}) ---
        documentation for the web admin portal (\texttt{admin.macsync.org}):
        portal overview, login and navigation, RBAC (built-in and custom roles,
        managing role per event, System Admin), creating/editing/deleting events
        and field reference, table/bus configuration, registered attendees and
        check-in columns, venue report PDF download, sending attendee
        notifications (including cooldown), user search and admin actions,
        warnings for destructive operations, and glossary terms (RBAC, managing
        role, venue report, etc.).
\end{itemize}
The manuals use consistent structure (document history, TOC, callouts for
notes/warnings/tips) so MES can hand them to members and staff at rollout.

\subsubsection*{Extra 2 --- Usability testing report}

We documented a structured usability evaluation in
\url{docs/Extras/UsabilityTesting/UsabilityTesting.tex} (dated April~6,
2026), covering both surfaces called out in scope: the React Native/Expo mobile
app for student attendees and the Next.js web admin dashboard for MES
organizers. The report states the plan (recruitment up to five participants per
group, on-campus sessions, think-aloud, screen/audio recording with consent,
tasks AT-1--AT-5 and OT-1--OT-5) and quantitative targets (e.g.\ task completion,
errors, SUS $\geq 70$).

The included results summary reports ten participants (five per
group) in a first round (March~2026 per the report's results template): mean SUS
76.5 (attendees) and 72.0 (organizers), both graded ``Good,''
with task completion rates and mean times tabulated per task. A usability issues
log records ten findings (e.g.\ post-sign-up email verification clarity, bus
sign-up discoverability, refund entry point labelling, admin event-form
validation feedback, notification wording, role-management navigation, back-stack
behaviour, analytics revenue labelling, payment error copy, loading states) with
recommended fixes and Nielsen severity ratings; the conclusion notes five
Severity~3 issues to address before a live event and a planned second test
round after fixes. The report also captures ethical safeguards (consent,
anonymized IDs, secure storage/deletion of recordings) and an informed-consent
appendix.

\section{Design Iteration (LO11 (PrototypeIterate))}

MacSync did not start as a single, self-contained product. Early in the capstone,
the work was split across three groups, each advancing different features
and assumptions in parallel. That structure helped
us explore ideas quickly, but it also made end-to-end behaviour hard to guarantee:
interfaces, schedules, and UX conventions could diverge, and MES stakeholders
still needed one dependable system they could hand to students and
volunteers for real large events. After weighing those integration risks against
the society's need for something shippable and supportable, we reset and
built MacSync from scratch as a standalone application---a dedicated mobile
experience for members and a web admin portal, backed by our own services, so the
platform could work on its own without depending on other teams' codebases
or other teams' release schedules. The change was driven by client-facing goals (predictable
registration, payments, check-in, and staff workflows) more than by internal
convenience: a standalone product meant fewer ``whose repo is this?'' problems
and let us own quality for the full stack.

From that baseline, iteration took the form of modeling, building, and
re-evaluating. We moved from informal sketches and dispersed prototypes to formal
artifacts (requirements, architecture, module interface specs, and guides) so
that each revision could be reviewed against MES scenarios---paid and free
events, seating, notifications, roles, and refunds---rather than against
fragile cross-repo glue. Implementation followed the same loop: early builds
surfaced gaps (missing or placeholder flows, unclear admin paths) that we
prioritized by severity and by what organizers need on event day.

Structured evaluation tightened that loop. Usability testing with
McMaster attendees and MES-style organizers (documented in our extras report)
showed that the product was broadly usable---both cohorts reached ``Good'' SUS
scores---but also highlighted concrete friction: for example, confusion after
sign-up about email verification, hard-to-find bus sign-up and refund entry
points on mobile, and weak validation and clarity in parts of the admin
dashboard (including notification wording and role-management navigation).
Those findings did not merely add bug reports; they explained \emph{why} users
stalled (unclear next step, labels that did not match mental models, missing
feedback after actions), which directed design and implementation work toward
clearer first-run experiences, more discoverable secondary tasks, and stronger
error and loading feedback for staff.

We also iterated in response to review and pilot use of the product.
Section~1 notes Revision~0 gaps that mattered to MES operations---for example,
members needing a working refund path and staff needing to edit events after
creation rather than living with placeholders. Addressing those items aligned the
implementation with how the client actually runs formal and semi-formal events,
where policies and venues change and payment exceptions happen.

In hindsight, the largest iteration was strategic---moving from a
multi-group feature split to a single standalone MacSync---because it set the
constraints under which all later prototypes made sense. Subsequent iterations
refined that product through usability testing, feedback on real workflows, and
course reviews---not only through merging code internally.

\section{Design Decisions (LO12)}

\subsection{Project challenge level (LO12)}

We classify this capstone as high challenge for a four-person, two-term
software project. The work spans a member mobile client, a
web admin portal, server APIs and data, third-party
payments (Stripe), notifications, role-based access control
across multiple MES-style organizers, and real event-day situations like registration spikes, seating logistics,
refunds, and check-in. The team
also produced a full documentation stack (planning, hazard analysis, design,
verification, user materials) while repivoting from a multi-group split
effort to a standalone product (Section~3). Together, these factors exceed a
typical single-surface CRUD app in technical breadth, risk, and coordination
load.

\subsection{Extras (LO12)}

The instructor-approved extras for this project are: \textbf{Extra~1} --- user
documentation (the MacSync Member Guide and the MacSync Staff \& Admin Guide
under \path{docs/Extras/UserManual/}); and \textbf{Extra~2} --- usability testing (plan,
metrics, tasks, and summarized results in
\path{docs/Extras/UsabilityTesting/UsabilityTesting.tex}). Section~2 summarizes
their contents; they constitute additional Rev~1 deliverables beyond the core
course set.

\subsection{Limitations}

Several limitations shaped what we designed and how thoroughly we could validate
it. Team size and time cap ambition: we prioritized MES-critical paths
(event lifecycle, tickets, payments, roles, notifications, basic reporting) over
nice-to-haves and deferred some non-functional depth (the VnV Report explicitly
records scope gaps such as webhook idempotency, reliability monitoring, and audit
depth). Dependence on third parties (Stripe, mobile push infrastructure,
hosting) limits control: we integrate rather than replicate payment processing,
which reduces PCI scope but ties behaviour to external APIs and policies.
Single-organization pilot reality means assumptions about traffic and
abuse are educated but not production-proven at massive scale. Testing
economy: usability samples are small (Nielsen-style); findings are directional
and were used to prioritize UX and validation fixes rather than to claim
statistical certainty.

These limitations pushed us toward clear module specs (MIS/MG) so builds
did not depend on word-of-mouth, Stripe instead of home-grown payments,
straightforward wording in VnV about what we did not fully test, and
UX changes informed by usability sessions.

\subsection{Assumptions}

We assumed McMaster-authenticated membership for general users
(\texttt{@mcmaster.ca}-based sign-up), networked clients at use time,
and good-faith student organizers coordinating with MES governance.
We assumed staff using the admin portal have access to a modern browser
and that MES can designate System Admins sparingly for elevated actions.
We assumed Stripe remains an acceptable processor for society ticket
sales and that notification delivery via standard mobile channels is sufficient
for event comms (with in-app history as backup). Design choices (email-bound
accounts, RBAC with a managing role per event, QR check-in, PDF venue report)
rest on these assumptions; breaking them would require different identity,
offline, or compliance models.

\subsection{Constraints}

Hard course constraints included fixed deadlines, required deliverable
types, TA/instructor rubrics, and reflection/trace requirements between Rev~0 and
Rev~1. Institutional constraints: student IP ownership policy; no
industry NDA driving hidden requirements. Technical constraints:
deploying a credible demo on approved stacks (React Native/Expo, Next.js,
service backend), keeping a clear line from requirements to tests in our docs, and
keeping payment handling out of custom card storage. Client constraints:
MES needs something volunteers can run for real events---so features like
refunds and editable events are not optional polish; they drove Rev~1
implementation priorities.

\subsection{How limitations, assumptions, and constraints influenced design}

The combination led to a dual-surface architecture: mobile-first for
attendees (discovery, purchase, ticket/QR in pocket) and web for staff (dense
tables, role management, reports). RBAC with per-event managing roles
balances flexibility for many clubs against a need to limit who can edit payouts,
attendee data, and notifications. Stripe was chosen to bound security
scope and leverage hosted checkout instead of storing payment credentials.
Formal design documents offset team turnover risk and TA expectations for
implementable specs. Where constraints forbade exhaustive production hardening,
we documented residual risk (hazard analysis, VnV scope statements)
instead of overstating coverage.

\subsection{Revision 0 $\rightarrow$ Revision 1 --- traceability of changes}

Table~\ref{TblRev0Rev1Trace} lists the substantive documentation,
verification, and implementation changes between Rev~0 and Rev~1 so an instructor
does not need to infer them from other sections. (Section~1 gives
feedback-by-source detail; here the focus is a single deliverable-level
change list.)

\begin{longtable}{@{} p{2.35cm} p{4.55cm} p{6.1cm} @{}}
\caption{Rev 0 $\rightarrow$ Rev 1 --- explicit change traceability (LO12)} \label{TblRev0Rev1Trace} \\
\toprule
\textbf{Deliverable} & \textbf{Rev 0 situation / gap} & \textbf{Rev 1 change (what we did)} \\
\midrule
\endfirsthead
\caption[]{\textit{Continued}} \\
\toprule
\textbf{Deliverable} & \textbf{Rev 0 situation / gap} & \textbf{Rev 1 change (what we did)} \\
\midrule
\endhead
\midrule
\multicolumn{3}{r}{\textit{Continued on next page}} \\
\endfoot
\bottomrule
\endlastfoot

Problem Statement & Objective and integrated-project context were too implicit; security/privacy, platform scope, integrations, and performance were thin; stakeholder tiers were overly rigid; goals were weakly measurable; security/privacy goal missing. & Added explicit main objective and integration context; expanded security/privacy, mobile+web scope, integrations, and high-traffic expectations; replaced rigid stakeholder hierarchy with a unified list; rewrote goals with measurable criteria and rationale; added Security \& Privacy goal; reflection appendix split Q1/Q2 per member and Q3 team-wide per rubric. \\

Development Plan & Technology section lacked tool-specific rationale; roles lacked technical ownership; workflow needed contributor-grade detail; schedule was shallow; POC plan needed clearer failure modes/context; formatting and GitHub link issues; reflection structure. & Expanded expected technology with rationale, trade-offs, and limits; added primary/secondary technical ownership across frontend, mobile, backend, DB, CI/CD; added PR/review norms, labels, quality gates, Definition of Done; added timetable table; strengthened POC plan (context, risks, scenario demo, failed POC definition); removed forced \texttt{flushleft} in communication plan; pointed GitHub link to repository root; reflection appendix Q1/Q2/Q3 structure aligned with rubric. \\

Software Requirements (SRS) & Rev~0 draft gaps: testability of some FR/NFR statements, cross-document naming relative to HA/design, and keeping requirement IDs aligned with VnV as the product matured. & Rev~1 AsciiDoc under \path{docs/SRS/}: clearer acceptance-style wording and measurable NFR targets; harmonized LEMS scope and integration language with HA and design; updated scenario/requirement coverage for refunds and admin event edit; feedback trace in Table~\ref{TblSRSFeedbackResponse}; explicit SRS$\leftrightarrow$HA narrative in Section~1. \\

Hazard Analysis & First draft; TA requested a system-boundary interaction diagram and clearer alignment/readability. & Added system-boundary figure and narrative alignment (LEMS application, database, external payment/notification dependencies); spelling/grammar polish (intro typo, out-of-scope list, FMEA part caption); traceability table in Reflect \& Trace (Section~1, Table~\ref{TblHazardAnalysisFeedbackResponse}). \\

Module Interface Spec (MIS) & Early MIS draft; symbols deferred to other docs; some module state underspecified for independent build; placeholders; cross-import noise. & Version~1.1: expanded in-document symbols/acronyms; sufficiency-for-build note per module; explicit state for payment (M1), sign-up (M3), access control (M5); session-timeout expectations internalized; appendix completed; Team~10 History Manager contribution linked (\url{https://github.com/OmarHassanAdelhamid/Five-of-a-Kind-capstone-project-/pull/328}); removed unused SRS label imports; TA review items incorporated per revision history. \\

Module Guide (MG) & Figures lacked narrative; caption typo; abbreviations depended on other docs. & Version~1.1: explanatory text before UI figures; fixed caption typo; MG defines its own abbreviations for consistency with MIS self-containment. \\

VnV Plan & Symbol table placeholder-heavy; template prompts remained; load/usability targets vague; unit-test detail waited on MIS. & Versions~1.1--1.2: expanded acronyms; removed \texttt{wss} instructor placeholders; concrete elevated-load definition tied to NFR cases; measurable usability criteria (tasks, timing, Likert for NFR3); appendix parameters; implementation verification populated with \texttt{UT-M*} cases for modules M1--M8, MIS mapping, SRS trace table. \\

VnV Report & Trace matrices scaled poorly; NFR IDs misaligned with plan; wording overstated completeness of load/usability; product gaps on refund and admin edit. & Landscape/wider matrices (no aggressive \texttt{resizebox}); NFR labels aligned to NFR1--NFR6 with crosswalk table; scope language clarified partial methods and gaps; implemented end-to-end refund request path and working admin event edit (backend + frontend). \\

Reflect \& Trace & Template-only or incomplete reflection on feedback and LO coverage. & Added structured feedback tables (PS, Dev Plan, Hazard Analysis), design-doc narrative (MIS/MG), VnV bullet responses, Extras (manuals + usability report), LO11 iteration narrative, and this LO12 section with explicit Rev~0$\rightarrow$~Rev~1 table. \\

Extras (beyond core set) & Not part of core Rev~0 package. & Delivered Member Guide, Staff \& Admin Guide, and Usability Testing report under \url{docs/Extras/UserManual/} and \url{docs/Extras/UsabilityTesting/} (Section~2). \\

\end{longtable}

\noindent\textit{SRS closing note.} The Software Requirements Specification remains the behavioural baseline cited by the
MIS/MG and VnV artifacts. Rev~1 verification (unit cases \texttt{UT-M*}, FR/NFR matrices, and scope honesty in the VnV Report)
traces execution and documentation back to that baseline; Table~\ref{TblSRSFeedbackResponse} and the SRS row of
Table~\ref{TblRev0Rev1Trace} document how Rev~1 SRS edits respond to feedback alongside the HA alignment described above.

\section{Economic Considerations (LO23)}

MacSync is not primarily a retail product: it is a society-owned event
platform for the McMaster Engineering Society (MES), funded through student fees
and club operations rather than through direct app sales. So the money story is
less about selling copies of the app and more about whether MES gets
value and what it costs to keep the system running.

Is there a ``market''? There is clear demand for the problem
MacSync solves. MES and affiliated clubs already run large events with hundreds
of attendees; today they stitch together forms, spreadsheets, social posts, and
ad hoc payments. A centralized registration, ticketing, seating, notification,
and check-in system reduces volunteer labour, cuts error and dispute risk around
payments and attendance, and gives students a single trustworthy place to act.
Comparable needs exist at other student societies, but our first
``market'' is MES-run programming at McMaster.

Marketing / adoption. For us, ``marketing'' really meant getting
people to trust and use the app, not big ad campaigns. Concretely, that includes: MES
endorsing the app as the official channel for listed events; club executives
being trained on the admin portal; listings in the MacSync mobile app on the
Apple App Store and Google Play Store; orientation-adjacent comms (e.g.\
society newsletters, Discord/Instagram pointers); and optional demo sessions for
organizers. Success is measured by active registrations and repeat use
per event, not by download counts alone.

Cost to produce a deployable version. A rough order-of-magnitude
for a version suitable for real MES events (excluding sunk student labour)
includes: cloud hosting for APIs and web admin (tens to low hundreds
of CAD per month depending on traffic and redundancy choices); domains
and certificates; Apple Developer and Google Play programme
fees for store distribution; Stripe transaction fees on ticket sales
(percentage-based, passed through or absorbed by policy); and admin
time for society volunteers (watching for errors, refunds, role changes). If we
priced our build time at industry salaries, the project would look expensive; as
a capstone, student hours stand in for a lot of that cash cost.

What would we charge? We would not charge a separate fee for
downloading MacSync. Revenue belongs to ticketed events: prices are set
per event, with Stripe handling card presentment. Any ``price'' for the platform
itself would be an internal MES line item (e.g.\ hosting subsidy from the
society), not something students pay for separately from tickets.

Units to break even. A classic break-even on app unit sales does
not apply. Break-even is better framed as: does event margin after
payment processing cover ongoing hosting and support? For a volunteer-run
society, even a small surplus on major formals, combined with avoidance of
manual reconciliation time, can justify fixed operating costs---provided incidents
(failed payments, refund disputes) stay manageable.

Attracting users if not ``sold.'' The practical approach is to tie the app
to events people already want: formals, pub nights, and similar. If those
events use MacSync for sign-up, students do not have to ``discover'' the app by
accident. Good guides (our extras), notifications that actually help, and quick
check-in at the door make it more likely someone will use the app again next time.

How many potential users exist? McMaster Engineering enrolls on
the order of thousands of undergraduate students; any of them may
attend MES-facing events. Typical large MES events draw hundreds of
attendees. Organizer-side users are smaller: tens to a few hundred volunteers and
executives across clubs with staff roles in the admin portal. That gives a
realistic picture of how many people might use MacSync for a first
deployment at McMaster, without pretending we are launching province-wide on day
one.

\section{Reflection on Project Management (LO24)}

We ran the project with a few habits that worked for us: at least one
full team meeting every week (and more when a deadline was close), and
written updates most other days---Discord messages, comments on GitHub
issues and PRs, and quick notes so people could read what changed when they had
time. We still used agendas and meeting minutes for the weekly call. GitHub kept
the work organized: issues and the board showed who was doing what, and pull
requests plus Actions helped catch broken builds before they hit \texttt{main}.
Live meetings were for bigger decisions and anything that was faster to sort out
by talking; everything else did not need another call. The rest of this section
compares what we did to the Development Plan, what worked, what did not, and
what we would change next time.

\subsection{How Does Your Project Management Compare to Your Development Plan}

Meetings and communication (largely followed). The Development Plan
originally described two standing team meetings per week; in practice
we standardized on at least one scheduled team meeting every week,
with additional syncs added ad hoc before major deadlines. Between
those calls, coordination was mostly in writing: Discord for quick
questions, GitHub for anything we wanted a record of on issues and PRs, and
shared docs when we were editing rubrics or specs together. That way we still
used agendas and meeting minutes from the chair, but nobody had
to be online twice a week unless they wanted to.
Between Rev~0 and the final deliverables we held seven internal
meetings in the documented contribution window, with full member
attendance (occasionally with late arrival or early departure). The plan's
12-hour response expectation was an aspiration for text and review
turnaround: it steered urgency on blockers but was not perfectly met during
exam and course-crunch periods.

Roles and workflow (followed in spirit with evolution). We kept the
named coordinating roles (leader, organizer, reviewer, notetaker/chair,
troubleshooter) and the technical ownership split (web primary/secondary,
mobile, backend, database, CI/CD). Issues were pre-created for milestones
then assigned in meetings so load balanced without relying on ad hoc
self-claiming alone. The workflow we documented---feature/\texttt{dev}/\texttt{main} branching,
small PRs, at least one review, stronger review for auth/payments/schema, and
a Definition of Done---matched how we aimed to run the repo.

Technology stack (aligned). We delivered on the planned ecosystem:
React Native / Expo for the member app, Next.js for the admin
dashboard, service/backend APIs with persisted data, Stripe for
payments, and GitHub Actions for CI-style checks on integration paths.
The plan's intent to stage via \texttt{dev} and release from \texttt{main} matches
our operating model even when course deadlines forced pragmatic shortcuts on
``production'' versus ``demo'' semantics.

Where we deviated from the written plan. The Development Plan assumed
weekly supervisor touchpoints during the tutorial slot; after the
strategic pivot to a standalone MacSync (see Section~3), we
consolidated supervisor engagement and held one focused
stakeholder meeting in the Rev~0-to-final window to align on approach rather
than weekly standing check-ins. We also missed the TA document
discussion session entirely (scheduling conflicts, religious observance, and
injury among members---as recorded in our final team contribution notes). Only
one of five members attended the single scheduled lecture in that period;
the rest relied on a teammate summary--- workable but not ideal for equal
exposure to instructor messaging. These problems came down to scheduling
and life outside the project, not to the team skipping our own meetings on
purpose.

\subsection{What Went Well?}

\begin{itemize}
  \item \textbf{Seeing what was actually in progress.} GitHub issues, milestones,
        and the Kanban columns (\texttt{Ready} $\rightarrow$ \texttt{In Progress}
        $\rightarrow$ \texttt{In Review} $\rightarrow$ \texttt{Done}) made it
        obvious when we were taking on too much---especially when course docs and
        VnV sat beside coding work.
  \item \textbf{A steady weekly rhythm.} Meeting once a week (sometimes more),
        plus short written updates on issues and PRs, meant problems did not sit
        hidden until the night before a deadline. We could move tasks around when
        someone got stuck.
  \item \textbf{Work spread across the team.} Commits and issue assignments stayed
        spread out enough that one person was never the only one who knew how the
        whole repo worked.
  \item \textbf{Pull request reviews.} We tried to review before merging, and to
        look harder at anything touching login, payments, or the database. That
        caught mistakes earlier than if we had only tested right before a demo.
  \item \textbf{Team charter expectations.} The charter asked for attendance,
        steady GitHub activity, and hitting dates we set in meetings. Those rules
        felt fair: when someone slipped, they said so early, and we still shipped
        on time.
\end{itemize}

\subsection{What Went Wrong?}

\begin{itemize}
  \item \textbf{We missed the TA document session.} None of us made it to the
        scheduled TA discussion, so we lost easy face-to-face feedback on drafts.
        We still went through rubrics and comments on our own, but skipping that
        meeting was a real miss.
  \item \textbf{Fewer supervisor check-ins after we changed direction.} The
        Development Plan assumed more regular contact with our supervisor; after
        we pivoted to a standalone MacSync we only held one focused meeting in the
        Rev~0-to-final window. That saved time but also meant less ongoing feedback
        from someone outside the team.
  \item \textbf{Only one person went to lecture.} Everyone else depended on a
        summary. If that one person missed a detail about a rubric change, the
        whole team could have been wrong together.
  \item \textbf{One person opened almost all the issues.} It kept labels and
        titles consistent, but when that person was busy, new work sometimes sat
        longer before it showed up on the board---we fixed some of that in
        meetings, but not through a second ``issue opener'' by default.
\end{itemize}

\subsection{What Would you Do Differently Next Time?}

\begin{itemize}
  \item \textbf{Book TA and tutorial time early.} Put TA sessions and important
        lectures on the calendar as soon as dates are known, and take turns being
        the person who goes in person so everyone gets a written summary the
        same day.
  \item \textbf{Short, regular supervisor check-ins after a big change.} After a
        pivot, a 15--20 minute chat every couple of weeks would have kept our
        supervisor in the loop without needing a long meeting every time.
  \item \textbf{Share who opens issues.} Use the same templates, but rotate who
        creates issues by milestone or area so the board never waits on one
        person.
  \item \textbf{Leave real time for docs and testing.} Course reports and VnV
        take the same people as the app; next time we would plan that overlap up
        front instead of assuming we could fit it in at the end.
\end{itemize}

\section{Ethics and Equity (LO18)}

Equity and universal design showed up in how we thought about who was
stressed when things went wrong, not only about abstract rules. Student
attendees mostly worry about wasting money, missing a seat on a bus, or standing
in a long line without a clear ticket---so we optimized for clear steps,
obvious errors with next actions, and tickets that work on a
phone. MES organizers worry about fairness between clubs, not overselling
capacity, and handling refunds without a spreadsheet war---so we invested in role
limits, audit-friendly flows in the admin tool, and documentation that does not
assume everyone is a developer. Thinking from those angles is how we tried to
practice empathy: we listened to usability testers describe where they got stuck
and read TA feedback about labeling and validation the same way an organizer or
first-year student might experience it.

Concrete steps toward equitable treatment. We tied sensitive attendee
data to permissions: staff see what they need for an event, not the
whole campus roster by default. The problem statement and hazard analysis already
flagged dietary and accessibility-related information and
privacy; design and verification work treated those categories as data
that deserves protection, not as optional notes. We aimed for a mobile-first
path for members because many students live on their phones and may not have a
laptop at the venue. User-facing manuals spelled out verification and
refund policy limitations (for example, who to contact if a refund stalls) so
members are not left guessing. Usability testing deliberately included
two kinds of stakeholders---attendee-style tasks and organizer-style
tasks---so we did not only optimize for our own team.

Universal design (what we actually built toward). Examples include:
clear navigation patterns in the mobile app (fixed tabs, consistent
language between docs and UI where we could); notification content that
works both as a push banner and inside the app for people who mute pushes;
large touch targets and simple flows in testing scenarios (buy ticket,
find QR, request refund); and honest error feedback called out in
usability findings (replacing vague ``something went wrong'' messages was flagged
as important because it helps everyone, not only expert users). Language in the
Member and Admin guides put warnings next to destructive actions (delete
user, delete event) so consequences are visible before someone clicks.

Where we fell short (honest misses). MacSync still assumes a valid
\texttt{@mcmaster.ca} account for ordinary registration. That matches MES's
world most of the time, but it is not equitable for guests, partners, or
students who do not use that email pattern without a manual workaround. We also
did not complete a full, formal WCAG-style accessibility audit with
screen-reader users in our capstone window; usability covered a small sample and
some checkpoints (contrast, clarity) but is not the same as certified compliance.
The problem statement mentioned richer dietary/accessibility matching
work than we fully automated; parts of that vision stayed at documentation or
process level rather than a polished product feature. Finally, heavy reliance on
smartphones helps many users but can exclude someone without a capable
device unless staff fall back to manual check-in---a gap to close with explicit
organizer playbooks, not only software.

Looking forward. If MacSync continues past the course, the next fairness
wins would be: documented offline/edge-case check-in procedures,
expanded assistive-technology testing, clearer pathways for
guests and others without a McMaster email if policy allows, and
measuring whether all major flows
remain usable after UI churn. Naming those gaps is part of the same ethical duty
as listing what we improved.

\section{Reflection on Capstone}

This capstone pulled together work we had only ever done in pieces during earlier
terms. The connections we sketched in our SRS-phase reflection
 still ring true at the end of the project: the
same courses kept paying off as we moved from requirements documents to a real
mobile app, admin web app, backend, and verification package.

\subsection{Which Courses Were Relevant}

\begin{itemize}
  \item \textbf{3RA3 (Software Requirements).} This course was the direct preparation for
        structuring requirements, naming stakeholders, and writing an
        SRS-style specification. It also introduced us to hazard thinking and the idea
        that verification should trace back to requirements---skills we reused in the
        Hazard Analysis, MIS-linked unit tests, and VnV artifacts instead of treating
        those documents as one-off essays.
  \item \textbf{3A04 (Software Design).} Design coursework helped us think in
        terms of modules, boundaries, and interfaces rather than a single giant
        program. That mindset matched how we separated attendee flows, admin flows,
        payments, and access control, and it made the Module Guide and MIS feel like
        natural extensions of how we already thought about the system.
  \item \textbf{2AA4 (Introduction to Software Development).} 2AA4 was our first
        serious experience building a multi-milestone group project: splitting
        features, iterating, and stacking each demo on the last. Capstone amplified that
        pattern (documentation sprints beside coding sprints), but the basic habit---agree
        on a thin vertical slice, then widen it---came from that course.
  \item \textbf{3S03 (Software Testing).} Testing coursework gave us a shared
        vocabulary for unit versus integration concerns, defect thinking, and
        why automated checks belong on pull requests. We leaned on those ideas when we
        argued for review gates and when we filled out implementation verification in
        the VnV Plan.
  \item \textbf{4HC3 (Human-Centred Interface Design).} HCI work pushed us to
        justify choices from a user and stakeholder perspective, not only from
        what is easiest to implement. That showed up in usability task design, in how
        we interpreted TA and tester feedback, and in the extras (Member and Admin
        guides written for non-developers).
\end{itemize}

\subsection{Knowledge/Skills Outside of Courses}

Even with those courses, we still had to learn on the job in a few areas
that the degree touches only lightly or not at all in one place:

\begin{itemize}
  \item \textbf{Stripe and payment flows.} Courses discussed security at a high
        level; integrating a live payment provider, webhooks, refunds, and
        sandbox testing required reading vendor docs and small experiments we
        owned ourselves.
  \item \textbf{React Native / Expo and modern front-end tooling.} We knew React
        ideas from design and web work, but mobile builds, store
        constraints, and Expo-specific debugging were mostly new in depth.
  \item \textbf{Production-shaped GitHub practice.} Branch protection,
        labels, project boards, and review etiquette at capstone scale went
        beyond what a single assignment usually enforces; we learned that by
        running the repo like a small team, not a homework fork.
  \item \textbf{Operations and deployment vocabulary.} Choosing hosting paths,
        environment secrets, and ``what breaks on demo day'' troubleshooting drew
        on tutorials and peers more than on any one lecture block.
  \item \textbf{Professional writing under rubrics.} CEAB-style reflection,
        trace tables, and instructor templates are a genre of their own; we got
        faster by reading past examples, TA comments, and each other's drafts.
\end{itemize}

In short, courses gave the skeleton (requirements, architecture,
testing mindset, user focus), and capstone forced us to put muscle on it
with stack-specific tools, integration risk, and documentation discipline we will
reuse in internships or future projects.

\end{document}